\documentclass{article}
\usepackage[T1]{fontenc}
\usepackage{lmodern}
\usepackage{graphicx}
\usepackage{amsmath,amssymb}
\usepackage[a4paper,margin=2cm]{geometry}
\usepackage[absolute,overlay]{textpos}
\setlength{\TPHorizModule}{1mm}
\setlength{\TPVertModule}{1mm}
\usepackage[indonesian]{babel}
\usepackage{hyperref}
\setlength{\emergencystretch}{2em}
% unit: Halaman 39

\begin{document}

% === Halaman 39 (python/native) ===

39

Contoh lain. Misalkan plainteks adalah 

CRYPTOGRAPHY AND DATA SECURITY 

Plainteks disusun menjadi 3 baris (k = 3) seperti di bawah ini: 

C 

 T 

  A 

   A     A     E       I 

 R   P  O R   P  Y  N  D  T  S  C   R   T 

   Y     G       H 

 D 

  A     U      Y 

maka cipherteksnya adalah 

CTAAAEIRPORPYNDTSCRTYGHDAUY 

2. Rail Fence Transposition Cipher

\end{document}
